% Generated by python/export_experiments.py; do not edit manually.
\begin{figure}[p]
\centering
\begin{tikzpicture}
\begin{axis}[diagnostic axis,width=.85\textwidth,height=7cm,xlabel={$a_q/a_{q,\max}$},ylabel={$\Delta R_{\mathrm{eq}}$ / m$\Omega$},title={Retain the spatial spread},legend columns=3,xmin=.1,xmax=1.02]
\addplot[reference samples] table[x=relative_iq,y=req_mOhm]{figures/data/experimental/outlier_2000_40_valid.dat};
\addlegendentry{40 $^\circ$C}
\addplot[measured samples] table[x=relative_iq,y=req_mOhm]{figures/data/experimental/outlier_2000_60_valid.dat};
\addlegendentry{60 $^\circ$C}
\addplot[only marks,derived quantity,mark=triangle*] table[x=relative_iq,y=req_mOhm]{figures/data/experimental/outlier_2000_80_valid.dat};
\addlegendentry{80 $^\circ$C}
\draw[guide quantity,guide pattern] (rel axis cs:.8152173913,0) -- (rel axis cs:.8152173913,1);
\end{axis}
\end{tikzpicture}
\caption{Each marker is one valid matched pair, projected onto relative measured q-current. The narrower spread near maximum q-current is an empirical feature of this dataset. The guide marks the initial 0.85 summary threshold; it does not define a universal identification region. No outlier removal is used.}
\label{fig:high-current}
\end{figure}
